% Language: LaTeX
% =====================================================================================================
%  ejemplos_de_notes.tex  --  Una muestra de cada tipo de nota definido en entornos.tex
%  Uso (apendices):  \input{secciones/6_apendices/ejemplos_de_notes}
%  Requiere: entornos.tex (float 'nota' y caja 'notaieee')
%
%  Tipos de nota:
%    notaieee{Titulo}     caja IEEE SIN numerar (solo diseno)
%    nota + notaieee      caja numerada "Nota X.Y": \caption y \label dentro del float 'nota'
%  Indice opcional:       \listofnota
% =====================================================================================================
\section{Ejemplos de notas}
\label{sec:ej-notas}

Este apartado muestra una muestra de cada tipo de nota definido en \texttt{entornos.tex}. Ambas variantes comparten el mismo diseño (caja azul estilo IEEE); la diferencia es que la segunda se numera, se puede referenciar y aparece en el índice de notas.

\subsection{Nota sin numerar}
\label{subsec:ej-nota-simple}

El entorno \texttt{notaieee} dibuja únicamente la caja y recibe el título como único argumento. Es adecuado para observaciones breves que no requieren referencia cruzada.

\begin{notaieee}{Nota IEEE sin numerar}
Contenido de la nota. La caja admite texto corrido, expresiones en línea como $E = mc^2$ y listas:
\begin{itemize}
  \item Primer punto de la nota.
  \item Segundo punto de la nota.
\end{itemize}
\end{notaieee}

\subsection{Nota numerada y referenciable}
\label{subsec:ej-nota-numerada}

Para numerar la nota, la caja \texttt{notaieee} se coloca dentro del float \texttt{nota}, que aporta el \texttt{\textbackslash caption} (numeración «Nota X.Y») y el \texttt{\textbackslash label} para las referencias cruzadas.

\begin{nota}
  \begin{notaieee}{Nota IEEE numerada}
  Esta nota se numera automáticamente y se registra en el índice de notas. Puede citarse en el texto con una referencia cruzada, como se hace a continuación.
  \end{notaieee}
  \caption{Ejemplo de nota numerada con \texttt{nota} y \texttt{notaieee}.}
  \label{nota:ej-numerada}
\end{nota}

Referencia cruzada a la numeración: Nota~\ref{nota:ej-numerada}.

% Indice de notas (opcional; lista todas las notas numeradas del documento):
% \listofnota